% Part: second-order-logic
% Chapter: metatheory
% Section: introduction

\documentclass[../../../include/open-logic-section]{subfiles}

\begin{document}

\olfileid[pt]{sol}{met}{int}

\olsection{Introdução}

A lógica de primeira ordem tem várias propriedades agradáveis. Sabemos que
ela não é decidível, mas ao menos é axiomatizável. Isto é, há sistemas de
prova corretos e completos para a lógica de primeira ordem, O.s.\@, eles
dão origem a !!a{derivability} relação~$\Proves$ com a propriedade de que,
para qualquer conjunto de !!{sentence}s~$\Gamma$ e qualquer
!!{sentence}~$!A$, $\Gamma \Entails !A$ sse $\Gamma \Proves !A$. Isso
significa, em particular, que as validades da lógica de primeira ordem são
!!{computably enumerable}. Existe uma função computável~$f\colon \Nat \to
\Sent[L]$ cujos valores de~$f$ são exatamente todas as !!{sentence}s válidas de
~$\Lang{L}$. Isso ocorre porque as !!{derivation}s podem ser enumeradas e
aquelas que !!{derive}m uma única~!!{sentence} são então associadas a essa
!!{sentence}. A lógica de segunda ordem é mais expressiva que a lógica de
primeira ordem e, portanto, em geral é mais complicado caracterizar suas
validades. De fato, mostraremos não apenas que a lógica de segunda ordem é
indecidível, mas também que suas validades nem sequer são !!{computably
  enumerable}. Isso significa que não pode haver um sistema de prova
correto e completo para a lógica de segunda ordem (embora existam sistemas
de prova corretos, porém incompletos, que de fato são importantes objetos
de pesquisa).

A lógica de primeira ordem também tem duas outras propriedades: ela é
compacta (se todo subconjunto finito de um conjunto~$\Gamma$ de
!!{sentence}s é satisfatível, então o próprio $\Gamma$ é satisfatível) e o
teorema de L\"owenheim--Skolem vale para ela (se $\Gamma$ tem um modelo
infinito, então tem !!a{denumerable} modelo). Ambos os resultados falham na
lógica de segunda ordem. Novamente, a razão é que a lógica de segunda ordem
pode expressar fatos sobre o tamanho de !!{domain}s que a lógica de primeira
ordem não pode expressar.

\end{document}
